Le stage de fin d'études constitue une étape charnière dans le parcours d'un étudiant en Licence. Il représente non seulement l'aboutissement de plusieurs années de formation théorique, mais aussi une immersion concrète dans le monde professionnel. C'est dans cette optique que ce projet a été entrepris au sein de GIAS Industries, une entreprise leader dans le secteur agroalimentaire en Tunisie.

Le choix de cette entreprise n'a pas été fortuit. Le secteur agroalimentaire, par sa complexité logistique et l'importance cruciale de son capital humain, offre un terrain d'étude particulièrement riche pour un futur diplômé en informatique de gestion. La motivation principale résidait dans le désir de confronter les acquis théoriques en gestion et en développement logiciel à des problématiques réelles de gestion des ressources humaines au sein d'une structure de grande envergure.

\vspace{1cm}

De nos jours, la transformation digitale est devenue un levier de performance incontournable pour les entreprises cherchant à optimiser leurs processus. Au sein de GIAS Industries, la gestion des stagiaires, bien que vitale pour le renouvellement des talents, faisait face à des défis liés à la dispersion des données et au manque d'automatisation. Ce projet s'inscrit donc dans une volonté de modernisation des services RH via la conception et la réalisation d'une plateforme web intégrée.

L'objectif est de digitaliser l'intégralité du cycle de vie du stagiaire : de la réception de la candidature à l'évaluation finale, en passant par le suivi des présences et la génération automatique des documents contractuels. Cette solution vise à libérer les responsables RH des tâches répétitives pour leur permettre de se concentrer sur des actions à plus forte valeur ajoutée.

\vspace{1cm}

Le présent rapport rend compte de ce travail et s'articule autour de trois chapitres principaux :

\begin{itemize}
    \item \textbf{Le premier chapitre} est dédié à la présentation de l'entreprise d'accueil, GIAS Industries, son secteur d'activité, son organisation interne ainsi que le département RH où s'est déroulé le stage.
    \item \textbf{Le deuxième chapitre} détaille le déroulement effectif du stage. Il expose les différentes phases de réalisation du projet, les technologies adoptées (React, Node.js, Sequelize) et les modules fonctionnels développés.
    \item \textbf{Le troisième chapitre} propose une analyse réflexive sur l'expérience vécue. Il aborde les problèmes techniques rencontrés, les solutions apportées, ainsi qu'un bilan des apports du stage tant pour l'étudiant que pour l'entreprise.
\end{itemize}

Enfin, une conclusion générale viendra clore ce rapport en synthétisant les résultats obtenus et en ouvrant des perspectives d'évolution pour la plateforme.
